% Appendix: training infrastructure, replication protocol, statistical significance.
% Numbers from SLURM accounting (sacct, all v14e-* jobs) and code_release/tests/binomial_ci.py.
% Regenerate the statistics with:  cd code_release && python tests/binomial_ci.py

\section{Training Infrastructure and Statistical Protocol}
\label{app:compute}

\subsection{Hardware and software}

All models were trained on an HPE Cray EX system with NVIDIA GH200
Grace--Hopper nodes: four GPUs per node, 95\,GiB of usable HBM per GPU, compute
capability 9.0, 128 CPU cores per task, and a Slingshot-11 interconnect
(NCCL over the \texttt{hsn} interfaces). Training used PyTorch 2.10.0 with CUDA
12.8, distributed data-parallel across nodes, and LoRA adapters over a frozen
Wan2.1-T2V-1.3B backbone.

The Default model and every conditioning ablation (PCA-2, PCA-4, No-Act-Tokens,
No-AdaLN, No-Critic) were trained on 8 nodes = 32 GPUs at global batch 32 and
learning rate $4.95\times10^{-5}$. The batch-size ablations used 4 nodes = 16
GPUs (global batch 16, lr $3.5\times10^{-5}$) and 16 nodes = 64 GPUs (global
batch 64, lr $7.0\times10^{-5}$). Every variant trains from step 0 to step 5000.
The scheduler enforces a four-hour limit per job, so each run proceeds as a
short sequence of checkpoint-resumed segments; Table~\ref{tab:compute} reports
the total.

\begin{table}[t]
\centering
\small
\begin{tabular}{lrrrr}
\toprule
Variant & Nodes & GPUs & Wall-clock & GPU-h \\
\midrule
Default (PCA-8)   & 8  & 32 & 26.9\,h & 862 \\
Batch64           & 16 & 64 & 28.4\,h & 1819 \\
Batch16           & 4  & 16 & 27.4\,h & 439 \\
PCA-4             & 8  & 32 & 27.5\,h & 879 \\
PCA-2             & 8  & 32 & 26.9\,h & 860 \\
No-Act-Tokens     & 8  & 32 & 21.0\,h & 671 \\
No-AdaLN          & 8  & 32 & 26.7\,h & 854 \\
No-Critic         & 8  & 32 & 26.8\,h & 857 \\
\midrule
\multicolumn{4}{l}{Total} & $\approx$7{,}250 \\
\bottomrule
\end{tabular}
\caption{Measured training cost per reported model.}
\label{tab:compute}
\end{table}

\subsection{One run per variant, many rollouts per run}

Each reported model is a \textbf{single training run} at \texttt{seed:\ 0},
offset per rank so that data order and initialisation are deterministic. We
report no seed replicates. At roughly 850 GPU-hours per variant, a three-seed
replication of the eight-model table would cost on the order of 22{,}000
GPU-hours, which was not available to us; we therefore spent the budget on
breadth of evaluation instead of depth of replication.

Concretely, we evaluate each trained model over \textbf{256 directional
rollouts} spanning the action grid and scene pool, and every quality column in
the paper is a statistic over that population rather than a single measurement.
The uncertainty we can quantify is therefore evaluation-sample uncertainty --
how much a reported percentage would move if we drew a different set of
rollouts from the same trained model -- and it is quantified exactly, from
hundreds of independent inferences per cell. What we cannot quantify from these
data is training-seed uncertainty, and we do not claim to.

Training is not bit-reproducible past the first step, because the backward pass
uses non-deterministic kernels. Measured run-to-run loss agreement is
$\sim$$10^{-3}$, against a loss that moves by $\sim$0.09 between adjacent steps
of a single run, so the residual is well inside one run's own variation.

\subsection{Significance}

The quality columns are proportions over known denominators --- 256 directional
rollouts for geometric corruption, and an explicitly defined feature-valid,
non-static active population for scene relocation. Proportions over known
denominators admit exact binomial inference directly from the published counts,
with no additional training runs. We report Wilson 95\% intervals per cell and
two-sided two-proportion $z$-tests between our model and each other model
(Table~\ref{tab:signif}).

\begin{table}[t]
\centering
\small
\begin{tabular}{lcc}
\toprule
Model & Geometric corruption & Scene relocation \\
      & \% [95\% CI]         & \% [95\% CI] \\
\midrule
MinWM      & \phantom{0}1.2 [0.4, 3.4]    & \phantom{0}3.7 [1.8, 7.5] \\
Yume       & \phantom{0}8.6 [5.7, 12.7]   & 57.8 [51.4, 63.9] \\
\textbf{Ours} & \textbf{16.8} [12.7, 21.9] & \textbf{10.9} [7.5, 15.5] \\
Batch64    & 25.8 [20.8, 31.5]            & \phantom{0}8.9 [5.9, 13.2] \\
Astra      & 52.0 [45.8, 58.0]            & 66.5 [60.0, 72.5] \\
WorldCam   & 68.0 [62.0, 73.4]            & 87.5 [82.4, 91.3] \\
MatrixGame & 70.7 [64.9, 75.9]            & 95.4 [92.0, 97.4] \\
WorldPlay  & 80.1 [74.8, 84.5]            & 77.3 [71.6, 82.2] \\
\bottomrule
\end{tabular}
\caption{Wilson 95\% confidence intervals on the two quality columns. Lower is
better in both.}
\label{tab:signif}
\end{table}

Our model's advantage over every external baseline is significant on both
columns ($p < 10^{-9}$ against Astra, WorldCam, MatrixGame and WorldPlay on
both; $p = 5.3\times10^{-3}$ against Yume on geometric corruption, $p < 10^{-9}$
on relocation). Two results run the other way and we state them plainly. MinWM
is significantly \emph{better} than our model on both columns
($p = 6.3\times10^{-10}$ geometry, $p = 6.3\times10^{-3}$ relocation), which is
the near-static failure mode discussed in the main text: a model that barely
moves has little opportunity to corrupt geometry or leave the scene. The
comparison against the Batch64 ablation splits by column: Default is
significantly better on geometric corruption (16.8\% vs.\ 25.8\%,
$p = 1.3\times10^{-2}$), but the relocation gap (10.9\% vs.\ 8.9\%) is
\textbf{not} significant ($p = 0.46$) and we do not claim a difference there.

These intervals characterise the rollout population for a fixed trained model.
They should not be read as intervals over training seeds, and the per-rollout
dispersion drawn in the response curves (SEM bands) and wedge plots
(interquartile arcs) likewise describes variation across rollouts within a
single run, not across runs.
